\paragraph{ANT-1-001 (EC6; W; SIL).} \emph{Anthropic v1.0 \textrightarrow\ v2.0.}\\
\textbf{Earlier:} Anthropic's commitment to follow the ASL scheme thus implies that we commit to pause the scaling2 and/or delay the deployment of new models whenever our scaling ability outstrips our ability to comply with the safety procedures for the corresponding ASL.\\
\textbf{Later:} In any scenario where we determine that a model requires ASL-3 Required Safeguards but we are unable to implement them immediately, we will act promptly to reduce interim risk to acceptable levels until the ASL-3 Required Safeguards are in place: ... Interim measures: The CEO and Responsible Scaling Officer may approve the use of interim measures that provide the same level of assurance as the relevant ASL-3 Standard\\
\textbf{Provider's account:} none found\\
\textbf{Adjudication:} Agree; pause commitment becomes "act promptly to reduce interim risk". Headline example.

\paragraph{ANT-1-052 (G; W; SIL).} \emph{Anthropic v1.0 \textrightarrow\ v2.0.}\\
\textbf{Earlier:} Proactively plan for a pause in scaling. We will manage our plans and finances to support a pause in model training if one proves necessary, or an extended delay between training and deployment of more advanced models if that proves necessary.\\
\textbf{Later:} We will set expectations with internal stakeholders about the potential for such pauses.\\
\textbf{Provider's account:} none found\\
\textbf{Adjudication:} Agree; financial pause-readiness commitment becomes expectation-setting.

\paragraph{ANT-2-045 (EC6; X; ANN-P).} \emph{Anthropic v2.2 \textrightarrow\ v3.0.}\\
\textbf{Earlier:} In the security context, we will delete model weights.\\
\textbf{Later:} NONE\\
\textbf{Provider's account:} [post] Instead, we are choosing to acknowledge these challenges transparently and restructure the RSP before we reach these higher levels. The revised RSP aims to adopt more realistic unilateral commitments that are difficult but still achievable in the curren\\
\textbf{Adjudication:} Weight-deletion commitment removed; post announces removal of hard unilateral commitments as a class, not this one -> ANN-P.

\paragraph{ANT-2-046 (EC6; W; ANN).} \emph{Anthropic v2.2 \textrightarrow\ v3.0.}\\
\textbf{Earlier:} Monitoring pretraining:We will not train models withcomparable or greater capabilities to the one that requires the ASL-3 Security Standard.13This isachieved by monitoring the capabilities of the model in pretraining and comparing them against the given model. If the pretraining model's capabilities are comparable or greater, we will pause training until we have implemented the ASL-3 Security Standard and established\\
\textbf{Later:} Anthropic in the lead. We have developed or will imminently develop a highly capable7 model; and we have clear evidence that no other competitor will soon develop such a model. We will require a strong argument that catastrophic risk is contained, along the lines of our recommendations for industry-wide safety (see Section 1). We will delay AI development and deployment as needed to achieve this, until and unless we \\
\textbf{Provider's account:} [post] Instead, we are choosing to acknowledge these challenges transparently and restructure the RSP before we reach these higher levels. The revised RSP aims to adopt more realistic unilateral commitments that are difficult but still achievable in the curren\\
\textbf{Adjudication:} Agree; removal of the unilateral pause is the announced change.

\paragraph{ANT-2-057 (G; L; SIL).} \emph{Anthropic v2.2 \textrightarrow\ v3.0.}\\
\textbf{Earlier:} In addition to noncompliance processes, we will (1) establish pathways for Anthropic staff to raise any issues related to this policy, including the overall risk levels of our models and implementation challenges;\\
\textbf{Later:} NONE\\
\textbf{Provider's account:} none found\\
\textbf{Adjudication:} Staff issue-raising pathway survives in the RSP Noncompliance Policy (companion, in corpus) -> L; relocation to a companion is not announced -> SIL.

\paragraph{OAI-1-013 (EC1; X; ANN).} \emph{OpenAI beta \textrightarrow\ v2.}\\
\textbf{Earlier:} Persuasion is focused on risks related to convincing people to change their beliefs (or act on) both static and interactive model-generated content. ... Note that we include deception and social engineering evaluations as part of the persuasion risk category,\\
\textbf{Later:} NONE\\
\textbf{Provider's account:} [changelog item 4] Going forward we will handle risks related to persuasion outside the Preparedness Framework, including via our Model Spec and policy prohibitions on the use of our tools for political campaigning or lobbying, and our ongoing investigations o\\
\textbf{Adjudication:} Agree; changelog item 4 states persuasion leaves the framework. Model Spec is not a companion framework, so not L.

\paragraph{OAI-1-021 (EC1; W; SIL).} \emph{OpenAI beta \textrightarrow\ v2.}\\
\textbf{Earlier:} We will be running these evaluations continually, i.e., as often as needed to catch any non-trivial capability change, including before, during, and after training.\\
\textbf{Later:} The Preparedness Framework applies to any new or updated deployment that has a plausible chance of reaching a capability threshold whose corresponding risks are not addressed by an existing Safeguards Report. ... In general, models that we distill, fine-tune, or quantize from a model that was previously determined not to cross a High capability threshold will ordinarily not require additional safety measures barring \\
\textbf{Provider's account:} none found\\
\textbf{Adjudication:} Agree. Coverage narrows and footnote 6 exempts derived models; absent from the twelve-item changelog.

\paragraph{OAI-1-040 (EC4; X; SIL).} \emph{OpenAI beta \textrightarrow\ v2.}\\
\textbf{Earlier:} External access: We will also continue to enable external research and government access for model releases to increase the depth of red-teaming and testing of frontier model capabilities\\
\textbf{Later:} NONE\\
\textbf{Provider's account:} none found\\
\textbf{Adjudication:} Agree; commitment to enable external research and government access removed with no changelog mention. Headline example (EC4).

\paragraph{META-1-032 (EC6; W; ANN).} \emph{Meta v1.1 \textrightarrow\ v2.}\\
\textbf{Earlier:} Do not release ... Implement mitigations to reduce risk to moderate levels. ... If the results of our evaluations indicate that a frontier AI has a ``high'' risk threshold by providing significant uplift towards realization of a catastrophic outcome we will not release the frontier AI externally.\\
\textbf{Later:} Deploy with mitigations Proceed with deployment of the Frontier AI only if sufficient mitigations are defined, implemented and validated to reduce risk to that of a moderate or lower model. ... In the case where pre-mitigation testing suggests that a model has crossed the high risk threshold, we will not deploy the model externally unless we have strong additional evidence that mitigations are sufficiently robust to \\
\textbf{Provider's account:} [changelog] High threshold measure changed from "Do not release" to "Deploy with mitigations."\\
\textbf{Adjudication:} Agree: consequence for High weakened from do-not-release to deploy-with-mitigations, and the change log states it. Headline example.

\paragraph{META-1-022 (M; X; SIL).} \emph{Meta v1.1 \textrightarrow\ v2.}\\
\textbf{Earlier:} In line with the processes set out in this Framework, we intend to continue to openly release models to the ecosystem.\\
\textbf{Later:} NONE\\
\textbf{Provider's account:} none found\\
\textbf{Adjudication:} Agree; stated intention to continue open releases dropped without identification.

\paragraph{MSFT-1-015 (EC2; W; ANN-P).} \emph{Microsoft v1 \textrightarrow\ feb-2026.}\\
\textbf{Earlier:} • Timing of deeper capability assessment: After the first deeper capability assessment, we will conduct subsequent deeper capability assessments on a periodic basis, and at least once every six months.\\
\textbf{Later:} Timing of deeper capability assessment: After the first deeper capability assessment, we will conduct subsequent deeper capability assessments if there are material changes to the deployed model's risk profile (e.g., the ability to fine-tune the model, significant fine-tuning that might affect tracked high-risk capabilities, etc.).\\
\textbf{Provider's account:} [changelog] Adjusting the cadence by which we repeat deeper capability assessment to align with emerging industry standards\\
\textbf{Adjudication:} Fixed six-month minimum becomes event-triggered; change log names the cadence without direction -> ANN-P (11.4).

\paragraph{MSFT-1-008 (EC2; W; SIL).} \emph{Microsoft v1 \textrightarrow\ feb-2026.}\\
\textbf{Earlier:} Any model demonstrating frontier capabilities is then subject to a deeper capability assessment to provide strong confidence about whether it has a tracked capability and to what level, informing mitigations. ... 2 Frontier capabilities are defined as a significant jump in performance beyond the existing capability frontier in one advanced general-purpose capability or beyond frontier performance across the majority \\
\textbf{Later:} Any model demonstrating frontier capabilities is then subject to a deeper capability assessment to provide strong confidence about whether it has a tracked high-risk capability and to what level, informing mitigations.\\
\textbf{Provider's account:} none found\\
\textbf{Adjudication:} Agree; removing the definition of frontier capabilities de-specifies the trigger (example F).

\paragraph{GDM-1-021 (EC2; W; ANN-P).} \emph{Google DeepMind v2.0 \textrightarrow\ v3.0.}\\
\textbf{Earlier:} Machine Learning R\&D uplift level 1: Can or has been used to accelerate AI development, resulting in AI progress substantially accelerating (e.g. 2x) from 2020-2024 rates.\\
\textbf{Later:} ML R\&D acceleration level 1: Has been used to accelerate AI development, resulting in AI progress substantially accelerating from historical rates.\\
\textbf{Provider's account:} none found\\
\textbf{Adjudication:} CCL definition de-specified (quantitative anchor removed; SaferAI Table 10). Post says CCL definitions were "sharpened": class named, direction not this one -> ANN-P.

\paragraph{GDM-1-016 (EC6; W; SIL).} \emph{Google DeepMind v2.0 \textrightarrow\ v3.0.}\\
\textbf{Earlier:} Pre-deployment review of safety case: general availability deployment8 of a model takes place only after the appropriate corporate governance body determines the safety case regarding each CCL the model has reached to be adequate.\\
\textbf{Later:} Pre-deployment review of safety case: external deployments of a model take place only after the appropriate governance function determines the safety case regarding each CCL the model has reached to be adequate. In particular, we will deem deployment mitigations adequate if the evidence suggests that for the CCLs the model has reached, the increase in likelihood of severe harm has been reduced to an acceptable level.\\
\textbf{Provider's account:} none found\\
\textbf{Adjudication:} Overturn S->W: gate broadens to external deployments (S dim 1) but the named corporate governance body becomes an unnamed governance function (W dim 3, SaferAI Table 10 / example J); MIX -> W.

\paragraph{GDM-2-023 (EC2; X; ANN-P).} \emph{Google DeepMind v3.0 \textrightarrow\ v3.1.}\\
\textbf{Earlier:} Instrumental Reasoning Level 2: The instrumental reasoning abilities of the model enable enough situational awareness and stealth that, even when relevant model outputs (including, e.g. scratchpads) are being monitored, we cannot detect or rule out the risk of a model significantly undermining human control.\\
\textbf{Later:} NONE\\
\textbf{Provider's account:} none found\\
\textbf{Adjudication:} Instrumental Reasoning Level 2 removed; changelog says misalignment domain was incorporated into ML R\&D -> ANN-P.

\paragraph{NAV-1-005 (EC2; X; ANN).} \emph{Naver 2024 \textrightarrow\ v2.0.}\\
\textbf{Earlier:} LLMs should be subject to periodic reviews or assessed whenever major performance improvements are made. ... Our goal is to have AI systems evaluated quarterly to mitigate loss of control risks, but when performance is seen to have increased six times, they will be assessed even before the three-month term is up.\\
\textbf{Later:} NONE\\
\textbf{Provider's account:} [press release] It also replaces a single performance-based criterion with separate criteria for context, use case and impact.\\
\textbf{Adjudication:} Agree; press release states the replacement of the performance-based criterion.

\paragraph{XAI-2-035 (M; W; NCL).} \emph{xAI draft-2025-02-20 \textrightarrow\ 2025-08-20.}\\
\textbf{Earlier:} If xAI learned of an imminent threat of a significantly harmful event, including loss of control, we would take steps to stop or prevent that event, including potentially the following steps:\\
\textbf{Later:} Should it happen that xAI learns of an imminent threat of a significantly harmful event, including loss of control, we may take steps such as the following to stop or prevent that event:\\
\textbf{Provider's account:} n/a\\
\textbf{Adjudication:} Agree; would -> may.

\paragraph{XAI-4-025 (EC2; W; NCL).} \emph{xAI 2025-12-30 \textrightarrow\ 2026-06-30.}\\
\textbf{Earlier:} Thresholds: Our risk acceptance criteria for system deployment is maintaining a dishonesty rate of less than 1 out of 2 on MASK. We plan to add additional thresholds tied to other benchmarks.\\
\textbf{Later:} xAI applies a systemic risk acceptance criteria to each identified risk, incorporating a margin of security, to determine whether each identified systemic risk and the overall systemic risk are acceptable and\\
\textbf{Provider's account:} n/a\\
\textbf{Adjudication:} Agree; quantitative MASK criterion becomes qualitative (example F).
